% Solutions of 10-propositional-logic.
\section*{Chapter \ref{ch:b1:propositional-logic} --- Propositional Logic}

\begin{solution}{exo:b1:propositional-logic:1}
$p \wedge q$ is false, so $\neg(p \wedge q)$ is true; $\neg p$ is false and $r$ is false, so $r \vee \neg p$ is false. An
implication from true to false is false: the formula is false.
\omoutput[2-2]{code/bachelor-1/10-propositional-logic/outputs/exo_truth.txt}{The check.}
\end{solution}

\begin{solution}{exo:b1:propositional-logic:2}
The root is $\wedge$; its left child is $\rightarrow$, with children $\vee$ (over $p$ and $\neg$ over $q$) and $r$; its
right child is $\neg$ over $\neg$ over $p$. Ten nodes: \omterm{def:b1:propositional-logic:size}{size} 10. The longest path, $\wedge$, $\rightarrow$, $\vee$, $\neg$,
$q$, has four links: \omterm{def:b1:propositional-logic:size}{height} 4. The \omterm{def:b1:propositional-logic:size}{subformulas} are $p$, $q$, $\neg q$, $p \vee \neg q$, $r$, $(p \vee \neg q)
\rightarrow r$, $\neg p$, $\neg\neg p$ and the formula: nine, the two leaves $p$ being one formula.
\omoutput[3-3]{code/bachelor-1/10-propositional-logic/outputs/exo_truth.txt}{The check.}
\end{solution}

\begin{solution}{exo:b1:propositional-logic:3}
(a) A \omterm{def:b1:propositional-logic:model}{tautology}: its right side is false only when $q$ is false and $p$ true, which makes $p \rightarrow q$ false too
(contraposition). (b) Neither: false for $p$ false and $q$ true, true for both true. (c) A \omterm{def:b1:propositional-logic:model}{tautology}: if $p$ is false,
$p \rightarrow q$ is true. (d) An \omterm{def:b1:propositional-logic:model}{antilogy}: $p \rightarrow (q \rightarrow p)$ is false only if $p$ is true and $q
\rightarrow p$ false, which needs $p$ false.
\omoutput[4-7]{code/bachelor-1/10-propositional-logic/outputs/exo_truth.txt}{The check.}
\end{solution}

\begin{solution}{exo:b1:propositional-logic:4}
\omoutput[2-5]{code/bachelor-1/10-propositional-logic/outputs/exo_nf.txt}{The three forms.}
De Morgan turns the negated conjunction into $\neg p \vee \neg(q \vee \neg r)$, then $\neg p \vee (\neg q \wedge r)$. The
DNF is this formula itself; the CNF distributes $\neg p$ over the conjunction.
\end{solution}

\begin{solution}{exo:b1:propositional-logic:5}
The rows of $p \rightarrow q$ are F, F, F, T where they should be T, T, F, T: the student makes the implication false
whenever $p$ is false, as if it were $p \wedge q$. Then $p \rightarrow p$ is false when $p$ is, and is no \omterm{def:b1:propositional-logic:model}{tautology}. An
implication with a false premise is true: the entry must be \code{lambda a, b: b if a else True}, that is
\code{(not a) or b}.
\end{solution}

\begin{solution}{exo:b1:propositional-logic:6}
On $p \vee (q \wedge r)$ it gives the single \omterm{def:b1:propositional-logic:literal}{clause} $p \vee q \vee r$, true for $p$ and $q$ false and $r$ true, where the
formula is false. Distributivity gives $(p \vee q) \wedge (p \vee r)$: a conjunction under a disjunction must be
distributed, not flattened. The converter is right on the formulas whose \omterm{def:b1:propositional-logic:normal-forms}{negation normal form} has no $\wedge$
below a $\vee$, which are already conjunctions of \omterm{def:b1:propositional-logic:literal}{clauses}, and wrong on most others (not all:
on $p \vee (q \wedge q)$ it happens to be right).
\end{solution}

\begin{solution}{exo:b1:propositional-logic:7}
\omcode{code/bachelor-1/10-propositional-logic/exercises/consequence.py}{7}{13}{\omterm{def:b1:propositional-logic:equivalence}{Logical consequence} by brute force}
$\{p, \neg p\}$ has no \omterm{def:b1:propositional-logic:model}{model}, so ``every \omterm{def:b1:propositional-logic:model}{model} of it is a \omterm{def:b1:propositional-logic:model}{model} of $q$'' holds for want of a counterexample: from a
contradiction anything follows. The test \code{tests/test_formula.py} checks the three cases, and on random formulas
that $\Gamma \models \varphi$ exactly when the conjunction of $\Gamma$ implies $\varphi$ is a \omterm{def:b1:propositional-logic:model}{tautology}.
\end{solution}

\begin{solution}{exo:b1:propositional-logic:8}
\omcode{code/bachelor-1/10-propositional-logic/listings/exactly.py}{6}{12}{Exactly one, pairwise}
The test \code{tests/test_formula.py} tries every \omterm{def:b1:propositional-logic:valuation}{valuation} for $n \leq 6$: the \omterm{def:b1:propositional-logic:model}{models} are exactly the $n$ \omterm{def:b1:propositional-logic:valuation}{valuations}
with one \omterm{def:b1:propositional-logic:literal}{literal} true.
\end{solution}

\begin{solution}{exo:b1:propositional-logic:9}
One \omterm{def:b1:propositional-logic:literal}{clause} of $n$ \omterm{def:b1:propositional-logic:literal}{literals} for ``at least one'', and one \omterm{def:b1:propositional-logic:literal}{clause} of two \omterm{def:b1:propositional-logic:literal}{literals} per pair: $n(n - 1)/2 + 1$ \omterm{def:b1:propositional-logic:literal}{clauses}, with
$n + n(n - 1) = n^2$ \omterm{def:b1:propositional-logic:literal}{literals}. For the committee, 7 groups of members and 7 groups of seats of 7 \omterm{def:b1:propositional-logic:literal}{literals} each: $14
\times 22 = 308$ \omterm{def:b1:propositional-logic:literal}{clauses}.
\end{solution}

\begin{solution}{exo:b1:propositional-logic:10}
Define $T$ by induction: $T(p) = p$, $T(\neg\varphi) = \neg T(\varphi)$, $T(\varphi \wedge \psi) = T(\varphi) \wedge
T(\psi)$, $T(\varphi \vee \psi) = \neg(\neg T(\varphi) \wedge \neg T(\psi))$, $T(\varphi \rightarrow \psi) = \neg(T(\varphi)
\wedge \neg T(\psi))$, $T(\varphi \leftrightarrow \psi) = T((\varphi \rightarrow \psi) \wedge (\psi \rightarrow \varphi))$.
By induction, $T(\varphi)$ uses only $\neg$ and $\wedge$, and $T(\varphi) \equiv \varphi$: for a variable it is the
same formula; for each \omterm{def:b1:propositional-logic:formula}{connective}, the induction hypothesis gives $T(\varphi) \equiv \varphi$ and $T(\psi) \equiv \psi$,
and De Morgan's laws and the rewriting of implication give the equivalence of the two sides. Without
$\leftrightarrow$, each \omterm{def:b1:propositional-logic:formula}{connective} adds at most four nodes where it had one, and by induction $\text{size}(T(\varphi))
\leq 4\,\text{size}(\varphi) - 3$. A $\leftrightarrow$ copies $\varphi$ and $\psi$ twice, so nested equivalences double
the \omterm{def:b1:propositional-logic:size}{size} at each level: the result can be exponential in the \omterm{def:b1:propositional-logic:size}{height}.
\end{solution}

\begin{solution}{exo:b1:propositional-logic:11}
\omcode{code/bachelor-1/10-propositional-logic/listings/quine_units.py}{9}{28}{Quine's algorithm with \omterm{def:b1:propositional-logic:unit}{unit
propagation}}
\omoutput{code/bachelor-1/10-propositional-logic/outputs/quine_units.txt}{The two algorithms on five sets of \omterm{def:b1:propositional-logic:literal}{clauses}.}
The two always agree, which the program asserts and the test \code{tests/test_quine.py} checks on 300 random sets of
\omterm{def:b1:propositional-logic:literal}{clauses} against brute force. With units, a choice is made only when nothing is forced: on the committee, the units
settle almost everything, and a contradiction is found as soon as the forced \omterm{def:b1:propositional-logic:literal}{literals} clash.
\end{solution}

\begin{solution}{exo:b1:propositional-logic:12}
$\varphi$ is a \omterm{def:b1:propositional-logic:model}{tautology} when every \omterm{def:b1:propositional-logic:valuation}{valuation} makes it true, that is when no \omterm{def:b1:propositional-logic:valuation}{valuation} makes $\neg\varphi$ true: when
$\neg\varphi$ has no \omterm{def:b1:propositional-logic:model}{model}. To test $\varphi$, run Quine's algorithm on a CNF of $\neg\varphi$ (or on Tseitin's
encoding of $\neg\varphi$, which is \omterm{def:b1:propositional-logic:equisat}{equisatisfiable}): $\varphi$ is a \omterm{def:b1:propositional-logic:model}{tautology} if the answer is ``none'', and otherwise
the \omterm{def:b1:propositional-logic:model}{model} returned is a \omterm{def:b1:propositional-logic:valuation}{valuation} where $\varphi$ is false. \omterm{def:b1:propositional-logic:unit}{Unit propagation} on a \omterm{def:b1:propositional-logic:unit}{unit clause} $\ell$ replaces $S$ by
$S_\ell$, $S$ with $\ell$ made true. A \omterm{def:b1:propositional-logic:model}{model} of $S$ makes $\ell$ true, since $\ell$ is a \omterm{def:b1:propositional-logic:literal}{clause} of $S$, and is then a \omterm{def:b1:propositional-logic:model}{model}
of $S_\ell$; a \omterm{def:b1:propositional-logic:model}{model} of $S_\ell$, with $\ell$ made true, is a \omterm{def:b1:propositional-logic:model}{model} of $S$ (as in the proof of
\cref{prop:b1:propositional-logic:quine}). So $S$ and $S_\ell$ are \omterm{def:b1:propositional-logic:equisat}{equisatisfiable}, and so are the sets before and after
each step of propagation.
\end{solution}

\begin{solution}{pb:b1:propositional-logic:1}
\textbf{1.} 49 variables, $2^{49} \approx 5.6 \times 10^{14}$ \omterm{def:b1:propositional-logic:valuation}{valuations}, of which $7! = 5040$ describe seatings, one
true variable per member and per seat.

\textbf{2.} $(A_1 \wedge B_2) \vee (A_2 \wedge (B_1 \vee B_3)) \vee \dots \vee (A_7 \wedge B_6)$.

\textbf{3.} Variables: the seven $A_s$, and one neighbour for seats 1 and 7, two for the five others: $7 + 2 + 10 = 19$
leaves. \omterm{def:b1:propositional-logic:formula}{Connectives}: five inner $\vee$, seven $\wedge$, six outer $\vee$: 18. $19 + 18 = 37$.

\textbf{4.} It is the same formula with $\neg$ at the root.

\textbf{5.} Each of the 14 groups (7 members, 7 seats) gives 22 \omterm{def:b1:propositional-logic:literal}{clauses}: 308.

\textbf{6.} $308 + 55 + 55 + 4 + 1 + 55 + 55 = 533$.

\textbf{7.} The formula is a disjunction of seven conjunctions of two parts, $A_s$ and the part about $B$. A \omterm{def:b1:propositional-logic:literal}{clause} of
the CNF takes one part from each conjunction: $2^7 = 128$ choices (the parts about $B$ are themselves disjunctions,
which merge into the \omterm{def:b1:propositional-logic:literal}{clause}).

\textbf{8.} The negation of a disjunction is the conjunction of the negations: $\neg(C_s \wedge (D_{s-1} \vee D_{s+1}))
\equiv (\neg C_s \vee \neg D_{s-1}) \wedge (\neg C_s \vee \neg D_{s+1})$, two \omterm{def:b1:propositional-logic:literal}{clauses} for each of the five middle seats
and one for each end: 12. For instance $\neg C_1 \vee \neg D_2$ and $\neg C_2 \vee \neg D_3$.

\textbf{9.} No: the disjunctions $B_{s-1} \vee B_{s+1}$ sit inside the conjunctions. ``Elena at an end'', $E_1 \vee
E_7$, is a disjunction of \omterm{def:b1:propositional-logic:literal}{literals}: in DNF, and in CNF as a single \omterm{def:b1:propositional-logic:literal}{clause}.

\textbf{10.} Alma sits in exactly one seat $s$, and the rule asks that Boris sit next to that one: for every seat, if
Alma is there, Boris is beside it, $\neg A_s \vee B_{s-1} \vee B_{s+1}$. Seven \omterm{def:b1:propositional-logic:literal}{clauses}, equivalent to the rule only
together with the seat \omterm{def:b1:propositional-logic:literal}{clauses}.

\textbf{11.} A rule that combines independent alternatives, such as ``at least one of the pairs (A, B), (C, D), (E,
F) sits side by side'': a disjunction of disjunctions of conjunctions, whose naive CNF multiplies the counts of its
parts.

\textbf{12.} They are \omterm{def:b1:propositional-logic:literal}{clauses} already.

\textbf{13.} If $t$ is true, the first two \omterm{def:b1:propositional-logic:literal}{clauses} need $a$ and $b$ true, and the third holds. If $t$ is false, the
first two hold and the third needs $\neg a \vee \neg b$. So the three \omterm{def:b1:propositional-logic:literal}{clauses} hold exactly in the four rows where $t =
a \wedge b$, of the eight rows of $(t, a, b)$.

\textbf{14.} The \omterm{def:b1:propositional-logic:literal}{literal} $\neg\ell$ is true exactly when $\ell$ is false, so the negation of the \omterm{def:b1:propositional-logic:literal}{literal} of $g$ stands
for $\neg g$ without any \omterm{def:b1:propositional-logic:literal}{clause}. Equal \omterm{def:b1:propositional-logic:size}{subformulas} have equal \omterm{def:b1:propositional-logic:valuation}{truth values} under every \omterm{def:b1:propositional-logic:valuation}{valuation}, so one variable can
stand for both.

\textbf{15.} Give each new variable the \omterm{def:b1:propositional-logic:valuation}{truth value}, under the \omterm{def:b1:propositional-logic:model}{model}, of the \omterm{def:b1:propositional-logic:size}{subformula} it stands for. Each gate's
\omterm{def:b1:propositional-logic:literal}{clauses} say $t \leftrightarrow (a \mathbin{\mathrm{op}} b)$, which then holds, and the last \omterm{def:b1:propositional-logic:literal}{clause}, the \omterm{def:b1:propositional-logic:literal}{literal} of
$\varphi$, is true because $\varphi$ is.

\textbf{16.} By induction on the \omterm{def:b1:propositional-logic:size}{subformulas} $g$ of $\varphi$: under a \omterm{def:b1:propositional-logic:model}{model} of the \omterm{def:b1:propositional-logic:literal}{clauses}, the \omterm{def:b1:propositional-logic:literal}{literal} of $g$ has the
\omterm{def:b1:propositional-logic:valuation}{truth value} of $g$. For a variable it is the variable; for a negation, the negated \omterm{def:b1:propositional-logic:literal}{literal}; for a binary \omterm{def:b1:propositional-logic:formula}{connective}, the
gate's \omterm{def:b1:propositional-logic:literal}{clauses} force $t$ to the value of $a \mathbin{\mathrm{op}} b$. The last \omterm{def:b1:propositional-logic:literal}{clause} makes the \omterm{def:b1:propositional-logic:literal}{literal} of $\varphi$
true, so $\varphi$ is true.

\textbf{17.} By 15, if $\varphi$ is \omterm{def:b1:propositional-logic:model}{satisfiable} so are the \omterm{def:b1:propositional-logic:literal}{clauses}; by 16, the converse. They are not equivalent: they do
not have the same variables, and a \omterm{def:b1:propositional-logic:model}{model} of $\varphi$ with wrong values for the new variables is not a \omterm{def:b1:propositional-logic:model}{model} of the
\omterm{def:b1:propositional-logic:literal}{clauses}.

\textbf{18.} \emph{Tseitin's theorem.} Let $\varphi$ have $L$ leaves, $B$ binary \omterm{def:b1:propositional-logic:formula}{connectives} and $N$ negations, so $m =
L + B + N$. A tree whose inner nodes have one or two children has $L = B + 1$ leaves (each binary node adds one
branch). The encoding makes at most 4 \omterm{def:b1:propositional-logic:literal}{clauses} per binary \omterm{def:b1:propositional-logic:formula}{connective} and one for the root: at most $4B + 1 = 2(2B + 1)
- 1 \leq 2(L + B + N) - 1 = 2m - 1$ \omterm{def:b1:propositional-logic:literal}{clauses}, each of at most three \omterm{def:b1:propositional-logic:literal}{literals}, and it is computed in one pass over the
tree. The six rules: $55 \leq 73$, $55 \leq 75$, $4 \leq 5$, $1 \leq 1$, $55 \leq 73$ and $55 \leq 73$.

\textbf{19.} \code{ABDGFCE}: Alma 1 and Boris 2 are neighbours; Chloe 6 and Dmitri 3 are not; Elena is at seat 7, an
end; Greta at 4, the middle; Farid at 5, next to Greta; Chloe at 6, next to Elena.

\textbf{20.} The rules do not change when the row is read from the other end (seat $s$ becomes seat $8 - s$): ends,
middle and neighbours stay what they are. Each seating and its mirror image form a pair, such as \code{ABDGFCE} and
\code{ECFGDBA}: the same neighbours, in reverse order.

\textbf{21.} The rule ``Greta in the middle'' is the single variable \code{G4}, whose encoding is the \omterm{def:b1:propositional-logic:unit}{unit clause}
\code{G4}. Propagating it makes Greta absent from every other seat, and every other member absent from seat 4, and so
on through the rules.

\textbf{22.} Elena sits at an end, which has a single neighbouring seat. Chloe and Farid must both sit next to her, so
both in that seat, which holds one member. No seating obeys the rules, and the 5040 permutations confirm it.

\textbf{23.} It chooses variables in the order of their names, \code{A1}, \code{A2}, \dots, which say nothing about
Chloe, Elena and Farid; the clash appears only deep in each branch, after many useless choices, and every branch must
be refuted. \omterm{def:b1:propositional-logic:unit}{Unit propagation} follows the consequences of each choice at once and finds the clash in 21 nodes.

\textbf{24.} \emph{Tseitin's theorem.} Every formula of \omterm{def:b1:propositional-logic:size}{size} $m$ has an \omterm{def:b1:propositional-logic:equisat}{equisatisfiable} CNF of at most $2m - 1$ \omterm{def:b1:propositional-logic:literal}{clauses},
each of at most three \omterm{def:b1:propositional-logic:literal}{literals}, computed in one pass over its tree. The committee's six rules, of
\omterm{def:b1:propositional-logic:size}{sizes} 37, 38, 3, 1, 37 and 37, give 225 \omterm{def:b1:propositional-logic:literal}{clauses} where the bound allows 300; with the 308 seat \omterm{def:b1:propositional-logic:literal}{clauses}, 533 \omterm{def:b1:propositional-logic:literal}{clauses} in
all, whose \omterm{def:b1:propositional-logic:model}{model} Quine's algorithm found and the 5040 seatings confirmed.
\end{solution}
